\documentclass[11pt,a4paper]{article}
\usepackage[T1]{fontenc}
\usepackage[utf8]{inputenc}
\usepackage{lmodern,amsmath,amssymb}
\usepackage[margin=25mm]{geometry}
\usepackage{longtable,booktabs,array,calc}
\usepackage[hidelinks]{hyperref}
\hypersetup{pdftitle={The dimension ladder: from 0+0 to 1+3, and what the action constant does on each rung},pdfauthor={navstokgap research working notes}}
\newcommand{\tightlist}{\setlength{\itemsep}{0pt}\setlength{\parskip}{0pt}}
\setlength{\emergencystretch}{3em}
\setcounter{secnumdepth}{0}
\begin{document}
\hypertarget{the-dimension-ladder-from-00-to-13-and-what-the-action-constant-does-on-each-rung}{%
\section{The dimension ladder: from 0+0 to 1+3, and what the action
constant does on each
rung}\label{the-dimension-ladder-from-00-to-13-and-what-the-action-constant-does-on-each-rung}}

\textbf{Record of a discussion with the user, 2026-09-27, 02:50--03:15.}
The user proposed the frame: quantization is a consistency condition
that Newton's rework of the \emph{Principia} missed, so any supposedly
well-posed wave equation without a quantization carries a logical issue;
the ladder of dimensions runs from \(0+0\) through Newton and quantum
mechanics (\(1+0\), ``YM\(_0\)'') to \(1+3\); and the exponential
\(e^{-1/(2b_0\hbar g^2)}\) plays a distinct role. This note records the
resulting statements with their status. The two layers of §1 and the
table of §2 are the parts meant for the joint paper; §§3--4 are
readings, labelled as such.

\hypertarget{quantization-as-a-consistency-condition-two-layers}{%
\subsection{1. Quantization as a consistency condition: two
layers}\label{quantization-as-a-consistency-condition-two-layers}}

\textbf{Layer 1: every classical wave equation needs a universal action
constant to be thermodynamically consistent.} Theorem U of the
\href{necessity-unit-and-indeterminacy.md}{unit-and-indeterminacy note}
uses Maxwell's field, but its mechanism is general. A classical field
whose normal-mode frequencies are unbounded (every wave equation on a
continuum) has, in equilibrium at temperature \(T\), energy \(k_BT\) per
mode by equipartition, hence infinite energy density. A finite
equilibrium, which the second law and every thermometer presuppose,
needs the high modes to freeze out; that compares \(k_BT\) with a
frequency and so needs a constant with units of action. Two different
fields in mutual equilibrium must share their occupation law, or a
filter between them would run a perpetual motion of the second kind
(Kirchhoff's argument extended from colours to fields); so the constant
is common to all fields. Einstein applied Planck's constant to the
vibrations of solids in exactly this way
(\href{https://doi.org/10.1002/andp.19063270110}{Einstein 1907},
metadata). Stated as a conjecture-level extension of Theorem U, to be
proved in the same form: \emph{any classical field with unbounded mode
frequencies that admits a finite thermal equilibrium requires a
universal action constant, and mutual equilibrium makes it the same for
all fields.} A well-posed wave equation remains consistent as
deterministic mathematics; as physics it is inconsistent once it can
exchange energy with thermal matter. This is the kind of gap the
programme attributes to Newton's rework: the limits of the
\emph{Principia} are consistent, and the inconsistency appears when they
are combined with thermodynamics and physical records.

\textbf{Layer 2: the quantized theory must itself survive refinement.}
Once \(\hbar\) is present, the per-cell parameter of the lattice
refinement is \(t=\hbar g_{\rm cl}^2a^{4-D}\)
(\href{zero-spacing-any-action.md}{zero-spacing note}, §3):

\begin{itemize}
\tightlist
\item
  \(D<4\): \(t\to0\) and the refinement converges (superrenormalizable);
\item
  \(D=4\): marginal; the theory survives only if its coupling runs to
  zero at short distances. Yang--Mills does (asymptotic freedom).
  Continuum \(\phi^4_4\) is trivial, a free field
  (\href{https://doi.org/10.1007/BF01205659}{Aizenman 1982};
  \href{https://doi.org/10.1016/0550-3213(82)90088-8}{Fröhlich 1982};
  \href{https://doi.org/10.4007/annals.2021.194.1.3}{Aizenman and
  Duminil-Copin 2021}; metadata), and QED has a Landau pole; on this
  reading both classical equations are effective theories with a cutoff;
\item
  gravity: \(t=\hbar G/(c^3a^2)=(\ell_P/a)^2\) grows as \(a\to0\), the
  failure mode of the \(D>4\) row.
\end{itemize}

\textbf{What ``consistency'' means in each layer.} Layer 1 is a logical
inconsistency: classical fields, the second law and finite energy are
jointly contradictory without an action constant. Layer 2 is an
existence condition on the continuum quantum theory, decided by
estimates of the kind proved in the series/parallel and zero-spacing
notes. Neither makes a classical PDE self-contradictory; both say it
cannot be a complete physical theory by itself. Classical Yang--Mills in
\(3+1\) Minkowski space exists globally
(\href{https://doi.org/10.1007/BF01976040}{Eardley and Moncrief 1982},
metadata), and classical \(\phi^4_4\) is well posed, although its
continuum quantum theory is free: the large-quantum-number limit of that
quantum theory is a free classical field. The implication ``quantum
theory exists \(\Rightarrow\) its classical limit exists'' holds in the
sense of limits; the converse fails as mathematics. The Euclidean
formulation makes one version exact: the measure \(e^{-S/\hbar}\) is a
classical statistical field theory in \(D\) dimensions with \(\hbar\) as
temperature, so ``quantum YM\(_4\) exists'' (Osterwalder--Schrader) and
``classical statistical YM\(_4\) at temperature \(\hbar\) exists in the
continuum'' are one problem, while the classical deterministic PDE is a
different object.

\hypertarget{the-ladder-in-the-users-1n-labelling}{%
\subsection{\texorpdfstring{2. The ladder, in the user's \(1+n\)
labelling}{2. The ladder, in the user's 1+n labelling}}\label{the-ladder-in-the-users-1n-labelling}}

Write \(\lambda_D=\hbar g_{\rm cl}^2\) with
\([\lambda_D]={\rm length}^{D-4}\), \(D=1+n\), and hold the classical
coupling \(g_{\rm cl}\) fixed as \(\hbar\to0\).

\begin{longtable}[]{@{}
  >{\raggedright\arraybackslash}p{(\columnwidth - 8\tabcolsep) * \real{0.2000}}
  >{\raggedright\arraybackslash}p{(\columnwidth - 8\tabcolsep) * \real{0.2000}}
  >{\raggedright\arraybackslash}p{(\columnwidth - 8\tabcolsep) * \real{0.2000}}
  >{\raggedright\arraybackslash}p{(\columnwidth - 8\tabcolsep) * \real{0.2000}}
  >{\raggedright\arraybackslash}p{(\columnwidth - 8\tabcolsep) * \real{0.2000}}@{}}
\toprule\noalign{}
\begin{minipage}[b]{\linewidth}\raggedright
\(1+n\)
\end{minipage} & \begin{minipage}[b]{\linewidth}\raggedright
Theory
\end{minipage} & \begin{minipage}[b]{\linewidth}\raggedright
Scale from \(\hbar\), \(c\), \(g_{\rm cl}\)
\end{minipage} & \begin{minipage}[b]{\linewidth}\raggedright
\(\hbar\)-dependence
\end{minipage} & \begin{minipage}[b]{\linewidth}\raggedright
Status
\end{minipage} \\
\midrule\noalign{}
\endhead
\bottomrule\noalign{}
\endlastfoot
\(0+0\) & single integral \(\int dU\,e^{-S/\hbar}\) & no time, no
spectrum & series in \(\hbar\) plus \(e^{-c/\hbar}\) & integrals \\
\(1+0\) & Newton, QM (``YM\(_0\)'') & record floor on action; rotor
\(\hbar^2C_2/(2I)\) & \(\hbar\) (action), \(\hbar^2\) (rotor) & Planck
paper; rotor exact \\
\(1+1\) & YM\(_1\) & \(\hbar c\sqrt{\lambda_2}\); circle
\(\hbar c\lambda_2LC_2/2\) & \(\hbar^{3/2}\); \(\hbar^2L\) & exact; no
local particle \\
\(1+2\) & YM\(_2\) & \(C_3\hbar c\lambda_3\) & \(\hbar^2\) & \(C_3\)
open \\
\(1+3\) & YM\(_3\) & \(C_4\hbar c\Lambda\) &
\(e^{-1/(2b_0\hbar g_{\rm cl}^2)}\) & the Millennium problem \\
\end{longtable}

(The subscripts of the series/parallel and zero-spacing notes count
spacetime dimension \(D\); here YM\(_n\) is \(D=n+1\).)

\textbf{\(0+0\).} There is no time direction, so no Hamiltonian and no
gap. What remains is one integral. It is the building block of every
rung: the plaquette weight \(k_t(U)\) is a \(0+0\) theory, series moves
(convolution) and parallel moves (products) assemble higher dimensions
from it, and the elimination of one inserted Newtonian time (Proposition
2 of the \href{refinement-composition-and-limit.md}{refinement note}) is
the \(0+0\) piece of mechanics. A zero-dimensional integral such as
\(\int d\phi\,e^{-(\phi^2/2+g\phi^4)/\hbar}\) already has an asymptotic,
non-Borel-summable series in \(\hbar\) with ambiguities of order
\(e^{-c/\hbar}\). At large \(N\) the one-plaquette unitary matrix
integral, the building block of two-dimensional lattice Yang--Mills, has
a third-order phase transition
(\href{https://doi.org/10.1103/PhysRevD.21.446}{Gross and Witten 1980},
metadata): non-analyticity from a zero-dimensional integral in a limit
that plays the role of \(a\to0\).

\textbf{\(1+1\) on a circle is exactly a \(1+0\) rotor.} Eq. (10) of the
refinement note, \(E_R=\hbar c\lambda_2L\,C_2(R)/2\), equals
\(\hbar^2C_2(R)/(2I)\) with \(I=1/(c\,g_{\rm cl}^2L)\): a quantum rigid
rotor on the gauge group whose moment of inertia is classical (\(\hbar\)
cancels from \(I\)). On the infinite line there is no particle; the
string tension \(\sigma=\hbar c\lambda_2C_2/2\) between static charges
goes like \(\hbar^2\).

\textbf{The exponents interpolate.} The only energy formed from
\(\hbar\), \(c\) and a fixed classical coupling is
\(E\sim\hbar c\,\lambda_D^{1/(4-D)}\), which goes like
\(\hbar^{(5-D)/(4-D)}\): \(\hbar^{3/2}\) for \(D=2\), \(\hbar^2\) for
\(D=3\), and a diverging exponent as \(D\to4\). At \(D=4\) no power is
available, and the running coupling turns the divergent power into
\(e^{-1/(2b_0\hbar g_{\rm cl}^2)}\), as dimensional regularization sees
it (\((g^2)^{1/\varepsilon}\) with \(\varepsilon=4-D\) becomes an
exponential of \(-1/g^2\)). This is dimensional analysis; whether each
rung has a positive gap is the open content (\(C_3\), and the Millennium
problem).

\textbf{Classical Yang--Mills has no gap in any dimension}: its waves
are massless and carry arbitrarily small energy. On every rung the
large-quantum-number limit is the limit in which the gap disappears.

\textbf{Refinement runs in opposite directions in mechanics and in field
theory.} For \(D<4\), \(\hbar\to0\) and \(a\to0\) push the same
per-plaquette parameter \(t=\hbar g_{\rm cl}^2a^{4-D}\) to zero: per
plaquette, ``more classical'' and ``finer lattice'' are one direction,
while the global scale where a gap lives stays quantum. In mechanics,
refining the time cells makes each cell more quantum: the classical
polygon action per cell shrinks like \(\tau^3\), while the fluctuation
action per cell stays of order \(\hbar\).

\hypertarget{a-third-route-to-a-universal-action-gravity-with-a-least-length}{%
\subsection{3. A third route to a universal action: gravity with a least
length}\label{a-third-route-to-a-universal-action-gravity-with-a-least-length}}

Classical gravitation with finite light speed has no mass-independent
action (\(Gm^2/c\) depends on the mass). A fixed length \(\ell\)
supplies one: \(c^3\ell^2/G\). So \((G,c,\ell)\) is a third route to a
universal action, beside radiation thermodynamics (Theorem U) and the
Stoney unit \(k_e/c\), and the most Newtonian of the three: \(G\) is
Newton's constant, \(c\) is Rømer's (1676), and the missing ingredient
is exactly the least magnitude that the scholium closing Book I, Section
I refuses by citing Euclid X. The value needs
\(\ell=\ell_P\approx1.6\times10^{-35}\) m; Newton's interval of fits
(\(1/89000\) inch) gives \(c^3\Lambda^2/G\sim10^{22}\) J s, off by about
56 orders of magnitude. The route is logically Newtonian; its value is
outside Newton-age data.

\hypertarget{the-exponential-ladder-a-reading}{%
\subsection{4. The exponential ladder (a
reading)}\label{the-exponential-ladder-a-reading}}

The factor \(e^{-1/(2b_0\hbar g^2)}\) has four properties worth keeping
together.

\begin{itemize}
\tightlist
\item
  It vanishes to all orders in \(\hbar\), so no finite order of the loop
  expansion sees the gap.
\item
  In four dimensions the running converts it into a power of length,
  \(e^{-c/g^2(a)}=(a\Lambda)^{2b_0c}\): dimensional transmutation. This
  is why the large-field criterion of the zero-spacing note is a
  power-law condition, \(2b_0c>4\), and why small instantons go like
  \((a\Lambda)^{11N/3}\).
\item
  Every scale of pure YM\(_4\) then reads as a suppression with an
  ``action'': \(\Lambda^k\propto e^{-k/(2b_0g^2)}\). The gap behaves
  like a configuration of action \(1/(2b_0)\) in units of \(g^{-2}\),
  which is \(24\pi^2/33\approx7.2\) for \(SU(3)\), against
  \(8\pi^2\approx79\) for an instanton; the dislocation threshold is
  \(k=D=4\) times it. This is a reading, not a result.
\item
  It lies where perturbation theory is ambiguous: the series in \(g^2\)
  is not Borel summable, with ambiguities of order
  \(e^{-{\rm const}/g^2}\), powers of \(\Lambda\) ('t Hooft's
  renormalons). The gap is in the class of quantities the perturbative
  expansion cannot decide, which is why non-perturbative estimates of
  the H1 type are unavoidable.
\end{itemize}

The same exponential runs through the programme: Theorem 5's
\(e^{-\pi^2/(8\lambda_3a)}\) in the series/parallel note, the instanton
\(e^{-8\pi^2/g^2}\), Wien's tail \(e^{-h\nu/k_BT}\) (where quantization
enters Theorem U, through the lumps \(h\nu\)), and on Newton's side the
action as a phase \(e^{iS/\hbar}\) or, in Euclidean form,
\(e^{-S/\hbar}\). In each case the exponential is where an action
constant meets a threshold. The paper could organize the comparison by
what each exponential suppresses and what survives refinement once the
exponent is converted into a scale.

\hypertarget{consequence-for-state}{%
\subsection{5. Consequence for STATE}\label{consequence-for-state}}

No change to the queue. The table of §2 is proposed as the spine of the
joint paper's comparison section, and Layer 1 of §1 as a theorem to
prove in the form of Theorem U for a general wave field (with the
cross-field universality). The gravity route of §3 belongs with the
dimensional note when that note is next revised.
\end{document}
